\subsection{Optische Anforderungen}
Gestalterische Aspekte spielen ebenfalls eine Rolle, da die WordClock sowohl als technisches System als auch als Gestaltungselement eingesetzt wird.
Der wichtigste Aspekt ist die Lesbarkeit und gleichmäßige Ausleuchtung der Anzeige. 
Die Worte müssen aus typischen Betrachtungsabständen klar erkennbar und ausreichend stark beleuchtet sein, dass ein Kontrast zwischen aktiven und inaktiven Buchstaben entsteht.
Weiterhin ist eine homogene Lichtverteilung sicherzustellen, um Hotspots oder blenden beim Betrachten zu vermeiden.
Die Hauptanzeige soll in einem Raster mit gleicher Seitenlänge angeordnet sein.
Alle optischen Anforderungen sind in Tabelle~\ref{tab:Optische Anforderung} aufgelistet.

\begin{table}[H]
\caption{Optische Anforderungen der WordClock}
\centering
\begin{tabular}{p{0.5cm} p{15cm}}
\textbf{ID} & \textbf{Beschreibung} \\
\hline
O1 & Gute Lesbarkeit der Anzeige bei typischen Betrachtungsabständen \\
O2 & Ausreichende Helligkeit zur Erzeugung eines klaren Kontrasts zu inaktiven Elementen \\
O3 & Gleichmäßige Ausleuchtung der gesamten Anzeige ohne lokale Überhellungen \\
O4 & Vermeidung von Hotspots und Blendwirkungen bei der Betrachtung \\
O5 & Homogene Lichtverteilung über alle dargestellten Wörter \\
O6 & Anordnung der Hauptanzeige in einem gleichmäßigen Raster mit identischer Seitenlänge \\
\end{tabular}
\label{tab:Optische Anforderung}
\end{table}